\section{The universal domain}
\label{sec:ud}

\subsection{Setting and closure operations}

\begin{convention}[the universal domain]\label{conv:ud}
Fix algebraically closed fields $k\subseteq K$ such that $\trdeg(K/k)$ is infinite.
A \emph{tuple} is a finite sequence of elements of $K$; tuples are concatenated
freely ($ab$ denotes the concatenation of $a$ and $b$) and a finite set of tuples
is identified with the concatenation of its members.  For a set $A\subseteq K$ we
write $k(A)$ for the subfield generated by $k\cup A$.  A subfield $L\subseteq K$ is
\emph{small} if $L\supseteq k$ and $\trdeg(L/k)<\infty$.
\end{convention}

\begin{leannote}
In Lean the ambient field is a type \lean{K} with \lean{[Field K] [IsAlgClosed K]}
and a base \lean{k} with \lean{[Field k] [IsAlgClosed k] [Algebra k K]}; the
hypothesis on the transcendence degree is \lean{Cardinal.aleph0 ≤ Algebra.trdeg k K}.
Tuples are functions \lean{Fin n → K}.  In the reconstruction blueprint the ambient
field is only assumed to have $\trdeg\ge5$; \cref{rem:finite-trdeg} and
\cref{rem:psi-hyp}(c) explain how the present results are applied there (by passing to
an ambient extension of infinite transcendence degree).
\end{leannote}

\begin{definition}[closures, rank, independence]\label{def:closures}
Let $A,B,C\subseteq K$ be sets (in practice finite tuples or small fields).
\begin{enumerate}[nosep]
\item $\acl(A)$ is the relative algebraic closure of $k(A)$ in $K$ (\lean{AclGeom.racl}).
  It is an algebraically closed subfield of $K$ containing $k$.
\item $\dcl(A)$ is the \emph{perfect closure} of $k(A)$ in $K$: in characteristic
  $0$, $\dcl(A)=k(A)$; in characteristic $p>0$,
  $\dcl(A)=\{x\in K:\ x^{p^e}\in k(A)\text{ for some }e\ge 0\}$.
\item $\rk(A/C)=\trdeg\bigl(k(C,A)/k(C)\bigr)\in\mathbb N\cup\{\infty\}$; it is finite
  when $A$ is finite.  We write $\rk(A)=\rk(A/\emptyset)$.
\item $A\ind{C}B$ (``$A$ is independent from $B$ over $C$'') means
  $\rk(a/C\cup B)=\rk(a/C)$ for every finite tuple $a$ from $A$.
\item $A$ and $B$ are \emph{interalgebraic}, $A\interalg B$, if $\acl(A)=\acl(B)$; they
  are \emph{interalgebraic over $C$} if $\acl(CA)=\acl(CB)$; they are
  \emph{interdefinable over $C$} if $\dcl(CA)=\dcl(CB)$.
\end{enumerate}
\end{definition}

\begin{lemma}[rank calculus]\label{lem:rank}
For finite tuples $a,a',b,d,e$ and sets $B,C,D$ (independence of sets is as in
\cref{def:closures}(4): $B\ind{C}D$ iff $b\ind{C}D$ for every finite tuple $b$ from $B$):
\begin{enumerate}[label=\textup{(R\arabic*)},start=0]
\item\label{R:sub} (finite character, sub-tuples) $\rk(a/C)=\rk(a/C_0)$ for some finite
  $C_0\subseteq C$, and $a\subseteq\acl(C)$ iff $a\subseteq\acl(C_0)$ for some finite
  $C_0\subseteq C$; if $ab\ind{C}D$ then $a\ind{C}D$; $a\ind{C}D$ iff $a\ind{C}d$ for
  every finite tuple $d$ from $D$.
\item\label{R:add} (additivity) $\rk(ab/C)=\rk(a/C)+\rk(b/Ca)$.
\item\label{R:acl} $\rk(a/C)=\rk(a/\acl(C))$, and $\rk(a/C)=0$ iff $a\subseteq\acl(C)$.
\item\label{R:mono} if $C\subseteq C'$ then $\rk(a/C')\le\rk(a/C)$.
\item\label{R:sym} $a\ind{C}b$ iff $b\ind{C}a$.
\item\label{R:trans} $a\ind{C}bd$ iff $a\ind{C}b$ and $a\ind{Cb}d$.
\item\label{R:aclinv} if $\acl(Ca')=\acl(Ca)$ then $a\ind{C}B \iff a'\ind{C}B$; and
  $a\ind{C}B\iff a\ind{C}\acl(CB)\iff \acl(Ca)\ind{C}\acl(CB)$ (independence of the
  closures, in the sense of \cref{def:closures}(4)).
\item\label{R:meet} if $B\ind{C}D$ then $\acl(CB)\cap\acl(CD)=\acl(C)$.
\item\label{R:selfind} if $B\ind{C}B$ then $B\subseteq\acl(C)$.
\item\label{R:swap} if $b\ind{C}D$ and $a\ind{C}bD$ then $ab\ind{C}D$; if
  $a\ind{C}b$ and $ab\ind{C}D$ then $a\ind{C}bD$.
\item\label{R:basechange} if $e\ind{}abC$ then $a\ind{C}b\iff a\ind{Ce}b$, and
  $a\subseteq\acl(C)\iff a\subseteq\acl(Ce)$.
\end{enumerate}
\end{lemma}

\begin{proof}
\ref{R:sub}: $k(C,a)$ is generated over $k(C)$ by finitely many elements, and a
transcendence basis of $k(C,a)/k(C)$ together with the finitely many algebraic
relations among $a$ and that basis involve only finitely many elements $C_0$ of $C$;
hence $\rk(a/C)=\rk(a/C_0)$, and likewise an algebraic dependence of each $a_i$ on $C$
involves finitely many elements of $C$.  Sub-tuples: $\rk(a/CD)\le\rk(a/C)$ and
$\rk(b/CaD)\le\rk(b/Ca)$ by \ref{R:mono}, and the sum of the two is
$\rk(ab/CD)=\rk(ab/C)$ by \ref{R:add}, so both are equalities.  The last clause is
\ref{R:sub} for $\rk(a/CD)=\rk(a/CD_0)$ with $D_0$ finite.
\ref{R:add} is the tower law for transcendence degree
(\lean{trdeg\_add\_eq}, or \lean{lift\_trdeg\_add\_eq} for the universe-lifted form)
applied to $k(C)\subseteq k(C,a)\subseteq k(C,a,b)$.
\ref{R:acl}: $k(C,a)$ and $\acl(C)(a)$ have the same transcendence degree over $k(C)$
and over $\acl(C)$ respectively because $\acl(C)/k(C)$ is algebraic; the second clause
is the definition of algebraic closure.  \ref{R:mono} is monotonicity of transcendence
degree in the base.  \ref{R:sym}: by \ref{R:add}, $\rk(ab/C)=\rk(a/C)+\rk(b/Ca)
=\rk(b/C)+\rk(a/Cb)$, so $\rk(b/Ca)=\rk(b/C)$ iff $\rk(a/Cb)=\rk(a/C)$.
\ref{R:trans}: $\rk(a/Cbd)\le\rk(a/Cb)\le\rk(a/C)$ with equality throughout iff
both inequalities are equalities.  \ref{R:aclinv}: the first clause follows from \ref{R:acl} and \ref{R:add}
applied to $\rk(aa'/C\cdot)$ (for each finite $b$ from $B$, by \ref{R:sub}); for the
second, $\rk(a/CB)=\rk(a/\acl(CB))$ by \ref{R:acl}, and for a finite tuple $b'$ from
$\acl(CB)$ there is a finite $b$ from $B$ with $b'\subseteq\acl(Cb)$ (\ref{R:sub}), so
$a\ind{C}b'$ follows from $a\ind{C}b$ and the first clause.  \ref{R:meet}: let $d$ be
a finite tuple from $\acl(CB)\cap\acl(CD)$.  Choose finite $b$ from $B$ with
$d\subseteq\acl(Cb)$ (\ref{R:sub}); then $b\ind{C}D$ and \ref{R:aclinv} give $d\ind{C}D$,
so $\rk(d/C)=\rk(d/CD)=0$ and $d\subseteq\acl(C)$.  \ref{R:selfind}: for $b$ from $B$,
$\rk(b/C)=\rk(b/CB)=0$.
\ref{R:swap}: first claim: $\rk(ab/CD)=\rk(b/CD)+\rk(a/CbD)=\rk(b/C)+\rk(a/C)=
\rk(ab/C)$ using \ref{R:trans} for $a$.  Second claim: $\rk(a/CbD)=\rk(ab/CD)-\rk(b/CD)
\ge\rk(ab/C)-\rk(b/C)=\rk(a/C)$, and the reverse inequality is \ref{R:mono}.
\ref{R:basechange}: $e\ind{}abC$ gives $e\ind{C}ab$ (\ref{R:trans}), hence $ab\ind{C}e$
(\ref{R:sym}) and then $a\ind{Cb}e$, $a\ind{C}e$ by \ref{R:trans}; these say
$\rk(a/Ceb)=\rk(a/Cb)$ and $\rk(a/Ce)=\rk(a/C)$, which gives the first
equivalence; the second is the case $b=\emptyset$ with \ref{R:acl}.
\end{proof}

\begin{leannote}
The relative closure $\acl$ is \lean{AclGeom.racl}; the needed facts are
\lean{AclGeom.mem\_racl\_iff}, \lean{AclGeom.racl\_exchange},
\lean{AclGeom.mem\_racl\_adjoin\_base\_iff} (change of base field),
\lean{AclGeom.mem\_racl\_of\_mem\_racl\_insert} (the non-degenerate case $|B|=|D|=1$,
$b\notin\acl(D)$, of \ref{R:meet})
and \lean{AclGeom.algebraicIndependent\_iff\_forall\_notMem\_racl}.  Ranks of closed
subfields are available as \lean{AclGeom.RankEq} with the matroid bridge
\lean{AclGeom.rankEq\_iff\_eRk}; the tuple-level rank $\rk(a/C)$ is
\lean{Algebra.trdeg (IntermediateField.adjoin k C) (IntermediateField.adjoin k (C ∪ a))}
or, equivalently, the matroid rank of $a$ relative to $C$ in
\lean{AlgebraicIndependent.matroid k K}.  Lemma~\ref{lem:rank} should be proved once
for the matroid rank (submodularity and the closure axioms are available in Mathlib's
\lean{Matroid} API) and transported.
\end{leannote}

\subsection{Automorphisms}

Let $\Gamma=\Aut(K/k)$ and, for a subfield $L\subseteq K$, $\Gamma_L=\Aut(K/L)$.
Every $\sigma\in\Gamma$ acts on tuples, sets, subfields and polynomial ideals
coordinatewise.

\begin{lemma}[fresh elements]\label{lem:fresh}
For every small subfield $L$ and every $n$ there is a tuple $t\in K^n$ with
$\rk(t/L)=n$.
\end{lemma}
\begin{proof}
$\trdeg(K/L)=\trdeg(K/k)$ is infinite since $\trdeg(L/k)$ is finite
(\lean{trdeg\_add\_eq}); a transcendence basis of $K$ over $L$ has at least $n$
elements.  (\lean{AclGeom.fresh\_four\_of\_five\_le\_trdeg} is the special case used in
the reconstruction blueprint.)
\end{proof}

\begin{lemma}[extension of isomorphisms; conjugacy]\label{lem:conj}
Let $L\subseteq K$ be any subfield and let $a,a'\in K^n$ be tuples with the same
vanishing ideal over $L$:
$\{f\in L[X_1,\dots,X_n]: f(a)=0\}=\{f: f(a')=0\}$.  Then there is
$\sigma\in\Gamma_L$ with $\sigma(a)=a'$.
\end{lemma}
\begin{proof}
This is \lean{AclGeom.exists\_ambientAutomorphism\_of\_idealOf\_eq}, applied with base
field $L$: the two vanishing ideals agree, so $L(a)\cong L(a')$ over $L$ by
$a\mapsto a'$; both are finitely generated over $L$, so $\trdeg(K/L(a))=\trdeg(K/L(a'))$,
and the isomorphism extends along transcendence bases and then to the algebraic
closure $K$ (\lean{IsAlgClosure.equivOfEquiv}).
\end{proof}

\begin{lemma}[equivariance]\label{lem:equiv}
For $\sigma\in\Gamma$: $\sigma(\acl A)=\acl(\sigma A)$, $\sigma(\dcl A)=\dcl(\sigma A)$,
$\rk(\sigma a/\sigma C)=\rk(a/C)$, and $a\ind{C}B\iff\sigma a\ind{\sigma C}\sigma B$.
\end{lemma}
\begin{proof}
$\sigma$ is a field automorphism fixing $k$; all notions are defined from the field
structure and $k$.
\end{proof}

\begin{lemma}[finite orbits]\label{lem:orbit}
Let $L\subseteq K$ be any subfield containing $k$ and $x\in K$.  The orbit
$\Gamma_L\cdot x$ is finite iff $x\in\acl(L)$ (the relative algebraic closure of $L$ in
$K$).  If $x$ is transcendental over $L$, then $\Gamma_L\cdot x$ is the set of all
elements of $K$ transcendental over $L$.
\end{lemma}
\begin{proof}
If $x$ is algebraic over $L$, every $\sigma x$ is a root of the minimal polynomial of
$x$ over $L$, so the orbit is finite.  If $x$ is transcendental over $L$, then for any
$x'$ transcendental over $L$ the vanishing ideals of $x$ and $x'$ over $L$ are both
zero, and \cref{lem:conj} gives $\sigma\in\Gamma_L$ with $\sigma x=x'$; the elements
$x,x^2,x^3,\dots$ are pairwise distinct and transcendental over $L$, so the orbit is
infinite.
\end{proof}

\subsection{Loci, generic points, dimension}

\begin{definition}[loci]\label{def:locus}
Let $a\in K^n$ and let $L\subseteq K$ be a subfield.  The \emph{vanishing ideal} of $a$
over $L$ is $I_L(a)=\{f\in L[X]: f(a)=0\}$, a prime ideal of $L[X]=L[X_1,\dots,X_n]$
(\lean{AclGeom.idealOf}, \lean{AclGeom.idealOf\_isPrime}; equivalently
\lean{MvPolynomial.vanishingIdeal L \{a\}}).  If $L$ is algebraically closed, the
\emph{locus ideal} of $a$ over $L$ is the extended ideal
\[
  P_L(a)=I_L(a)\cdot K[X]\subseteq K[X].
\]
An ideal $I\subseteq K[X]$ is \emph{defined over} a subfield $F\subseteq K$ if
$I=(I\cap F[X])\cdot K[X]$.  A tuple $a'\in K^n$ is a \emph{generic point of a prime $P$
over an algebraically closed subfield $M$} if $P_M(a')=P$.
\end{definition}

\begin{lemma}[extension and contraction]\label{lem:extcontr}
Let $F\subseteq M\subseteq K$ be subfields and $J\subseteq F[X]$ an ideal.
\begin{enumerate}[nosep]
\item (change of coefficients) $(J\cdot K[X])\cap M[X]=J\cdot M[X]$ and
  $(J\cdot K[X])\cap F[X]=J$.
\item (dropping variables) If $J\subseteq M[Y,Z]$ is an ideal in two groups of variables,
  then $(J\cdot K[Y,Z])\cap K[Y]=(J\cap M[Y])\cdot K[Y]$; in particular an ideal of
  $K[Y,Z]$ defined over $M$ contracts to an ideal of $K[Y]$ defined over $M$.
\end{enumerate}
\end{lemma}
\begin{proof}
(1) Choose an $M$-basis $(\beta_i)$ of $K$ containing $1$.  Then $K[X]=\bigoplus_i M[X]\beta_i$
as $M[X]$-modules and $J\cdot K[X]=(J\cdot M[X])\cdot K[X]=\bigoplus_i (J\cdot M[X])\beta_i$;
intersecting with $M[X]$ (the summand of $\beta=1$) gives $J\cdot M[X]$.  The second
claim is the case $M=F$.  (\lean{MvPolynomial.ker\_map} and
\lean{Ideal.comap\_map\_eq\_self\_of\_faithfullyFlat} are the Mathlib ingredients:
$K[X]$ is faithfully flat over $M[X]$.)
(2) With the same basis, $K[Y,Z]=\bigoplus_i M[Y,Z]\beta_i$ and
$K[Y]=\bigoplus_i M[Y]\beta_i$; so $J\cdot K[Y,Z]=\bigoplus_i J\beta_i$ and an element
$\sum_i h_i\beta_i$ ($h_i\in J$) lies in $K[Y]$ iff every $h_i\in M[Y]$, i.e.\ iff every
$h_i\in J\cap M[Y]$.
\end{proof}

\begin{lemma}[locus ideals are prime]\label{lem:locusprime}
If $L$ is algebraically closed then $P_L(a)$ is a prime ideal of $K[X]$, and
$P_L(a)\cap L[X]=I_L(a)$.  More generally, for algebraically closed $L\subseteq M\subseteq K$
and a prime $I\subseteq L[X]$, the extension $I\cdot M[X]$ is prime.
\end{lemma}
\begin{proof}
\lean{AclGeom.isPrime\_map\_of\_isAlgClosed} (the extension of a prime ideal from an
algebraically closed field along any field extension is prime) and
\cref{lem:extcontr}.
\end{proof}

\begin{lemma}[existence of generic points]\label{lem:genexists}
Let $M$ be a small algebraically closed subfield and $P\subseteq K[X]$ a prime ideal
defined over $M$.  Then $P$ has a generic point over $M$.  Any two generic points of
$P$ over $M$ are conjugate under $\Gamma_M$, and every $\Gamma_M$-conjugate of a generic
point is generic.  Moreover (generic over a larger base implies generic over a smaller
one): if $F\subseteq M$ is algebraically closed, $P$ is defined over $F$, and
$P_M(a')=P$, then $P_F(a')=P$.
\end{lemma}
\begin{proof}
$Q=P\cap M[X]$ is prime and $P=Q\cdot K[X]$.  The domain $M[X]/Q$ is finitely
generated over $M$; its fraction field $E$ has finite transcendence degree $d$ over $M$.
Choose a transcendence basis $t_1,\dots,t_d$ of $E/M$ and elements $s_1,\dots,s_d\in K$
with $\rk(s/M)=d$ (\cref{lem:fresh}); the $M$-embedding $M(t)\to K$, $t_i\mapsto s_i$,
extends to $E$ since $E/M(t)$ is algebraic and $K$ is algebraically closed
(\lean{AclGeom.adjoinTranscendentalAlgHom},
\lean{AclGeom.exists\_extension\_of\_isAlgebraic}, or \lean{IsAlgClosed.lift}).  The image
$a'$ of the residue classes of $X_1,\dots,X_n$ satisfies $I_M(a')=Q$, i.e.\ $P_M(a')=P$.
If $a',a''$ are two generic points then $I_M(a')=I_M(a'')=Q$ and \cref{lem:conj} applies;
conversely $I_M(\sigma a')=\sigma(I_M(a'))=I_M(a')$ for $\sigma\in\Gamma_M$.
For the last clause, $P_F(a')\subseteq P_M(a')=P$ since $I_F(a')\subseteq I_M(a')$;
conversely $P\cap F[X]\subseteq P\cap M[X]=I_M(a')$ (\cref{lem:extcontr}), so
$P\cap F[X]\subseteq I_M(a')\cap F[X]=I_F(a')$ and
$P=(P\cap F[X])K[X]\subseteq I_F(a')K[X]=P_F(a')$.
\end{proof}

\begin{lemma}[dimension drop]\label{lem:dimdrop}
Let $M$ be any field and $Q\subseteq Q'$ prime ideals of $M[X]$.  Then
$\trdeg_M\bigl(\operatorname{Frac}(M[X]/Q')\bigr)\le\trdeg_M\bigl(\operatorname{Frac}(M[X]/Q)\bigr)$,
with strict inequality if $Q\subsetneq Q'$.
\end{lemma}
\begin{proof}
Write $A=M[X]/Q$, $A'=M[X]/Q'$, $\pi:A\to A'$ the surjection, $d'=\trdeg_M A'$ and
$d=\trdeg_M A$ (transcendence degrees of fraction fields).  Preimages of algebraically
independent elements are algebraically independent, so $d'\le d$.  Now let
$Q\subsetneq Q'$ and suppose $d'=d$.
Choose $t_1,\dots,t_d\in A'$ algebraically independent over $M$ and lifts $s_i\in A$;
then $s_1,\dots,s_d$ form a transcendence basis of $\operatorname{Frac}A$ over $M$.
Pick $f\in Q'\setminus Q$; its class $\bar f\in A$ is nonzero and $\pi(\bar f)=0$.
Since $\bar f$ is algebraic over $M(s)$ there is a nonzero $g\in M[S_1,\dots,S_d,T]$
with $g(s,\bar f)=0$; take $g$ of least $T$-degree.  Its constant term $g_0(S)$ is
nonzero: otherwise $g=T\cdot g_1$, and $\bar f g_1(s,\bar f)=0$ in the domain $A$ with
$\bar f\ne 0$ gives $g_1(s,\bar f)=0$ with smaller $T$-degree.  Applying $\pi$,
$0=\pi(g(s,\bar f))=g(t,0)=g_0(t)$, contradicting the algebraic independence of $t$.
\end{proof}

\begin{lemma}[dimension of a locus]\label{lem:dim}
Let $L\subseteq M$ be small algebraically closed subfields and $a\in K^n$.
\begin{enumerate}[nosep]
\item If $a'$ is a generic point of $P_L(a)$ over $M$ then $\rk(a'/M)=\rk(a/L)$.
\item (locus under an independent base extension) $P_M(a)=P_L(a)$ iff $a\ind{L}M$.
\item $P_L(a)=P_L(a')$ iff $a$ and $a'$ are $\Gamma_L$-conjugate.
\end{enumerate}
Consequently the \emph{dimension} $\dim P:=\rk(a'/M)$ of a prime $P\subseteq K[X]$
defined over a small algebraically closed $M$, computed with any generic point $a'$
over any such $M$, is well defined, and $\dim P_L(a)=\rk(a/L)$.
\end{lemma}
\begin{proof}
(1) By \cref{lem:locusprime,lem:extcontr}, $I_M(a')=P_L(a)\cap M[X]=I_L(a)\cdot M[X]$.
This says that $a'$ is a generic point over $M$ of the extension of the prime
$I_L(a)\subseteq L[X]$ to $M$; the theorem
\lean{AclGeom.trdeg\_adjoin\_of\_isGeneric} (with base $L$, extension $M$ and ideal
$I_L(a)$) gives $\rk(a'/M)=\trdeg_M M(a')=\trdeg_L L(a')$, and
\lean{AclGeom.isGeneric\_baseChange} gives $I_L(a')=I_L(a)$, so by \cref{lem:conj} $a'$
and $a$ are $\Gamma_L$-conjugate and $\trdeg_L L(a')=\rk(a/L)$.

(2) Always $P_L(a)\subseteq P_M(a)$ since $I_L(a)\subseteq I_M(a)$.  Both are prime
ideals; by (1) applied to a generic point $a'$ of $P_L(a)$ over $M$ and to $a$ itself,
$\trdeg_M\operatorname{Frac}(M[X]/(P_L(a)\cap M[X]))=\rk(a/L)$ and
$\trdeg_M\operatorname{Frac}(M[X]/(P_M(a)\cap M[X]))=\rk(a/M)$.  If $a\ind{L}M$ the
two numbers are equal and \cref{lem:dimdrop} forces $P_L(a)\cap M[X]=P_M(a)\cap M[X]$,
hence $P_L(a)=P_M(a)$ by \cref{lem:extcontr}.  Conversely if $P_L(a)=P_M(a)$ the two
numbers agree, i.e.\ $\rk(a/M)=\rk(a/L)$.

(3) is \cref{lem:conj} with the base $L$ together with \cref{lem:extcontr}
($P_L(a)=P_L(a')$ iff $I_L(a)=I_L(a')$).  The final statement follows from (1)
applied over $\acl(M\cup M')$ for two choices $M,M'$ of base.
\end{proof}

\begin{lemma}[conjugates independent from a given tuple]\label{lem:fresh-conj}
Let $M$ be a small algebraically closed subfield, $a$ a tuple and $w$ any tuple.  Then
there is $\sigma\in\Gamma_M$ with $\sigma a\ind{M}w$.  Consequently, if $\Phi$ is a
$\Gamma_M$-invariant property of tuples and $\Phi(a)$ holds, there is $a'$ with
$\Phi(a')$ and $a'\ind{M}w$.
\end{lemma}
\begin{proof}
Let $M'=\acl(M,w)$ and let $a'$ be a generic point of $P_M(a)$ over $M'$
(\cref{lem:genexists}).  By \cref{lem:dim}(1), $\rk(a'/M')=\rk(a/M)$; by the last
clause of \cref{lem:genexists}, $P_M(a')=P_M(a)$, so $a'=\sigma a$ for some
$\sigma\in\Gamma_M$ (\cref{lem:dim}(3)) and $\rk(a'/M)=\rk(a/M)=\rk(a'/M')$, i.e.\
$a'\ind{M}w$ (\ref{R:acl}).
\end{proof}

\begin{remark}[the infinite transcendence degree is essential]\label{rem:finite-trdeg}
The hypothesis $\trdeg(K/k)=\infty$ of \cref{conv:ud} is used in two ways.  Explicitly,
the proofs adjoin fresh generic points: the parameter tuple $e$ of \cref{thm:main} has
rank $5n$ (\cref{rem:params}), \cref{thm:constructible} adds a generic point $y_0$ and the
sampling tuple $t=(t_0,t_1)$, and the auxiliary lemmas (\cref{lem:genexists},
\cref{lem:fresh}, \cref{lem:fresh-conj}) adjoin further generic points in a nested way.
Implicitly, and more seriously, the definitions quantify over all \emph{small}
algebraically closed subfields: if $\trdeg(K/k)$ were finite, $K$ itself would be small,
a variety of positive dimension would have no generic point over $K$, and the pair
$(K,f)$ would be a representative of a germ vacuously, equivalent to every other
representative; the germ quotient of \cref{def:germ} would collapse.  So the results of
this document are \emph{not} claimed for ambient fields of finite transcendence degree,
and the parameter count $5n$ does not control the fields quantified over in the
definitions.  For an algebraically closed $K$ of finite transcendence degree over $k$
one passes to an algebraically closed extension $K'$ of infinite transcendence degree,
applies the theorems there, and then needs two separate descents: the parameter
descent of \cref{sec:descent}, which is proved here and moves the adjoined parameters
into $k$, and a bounded first-order (model completeness) argument, not proved here,
which returns existential witnesses from $K'$ to $K$ (\cref{rem:psi-hyp}(c)).
\end{remark}
